\color{unimain}{\subsection{Защита от неполнофазного режима  (ЗНР)}}
\color{black}

\begin{enumerate}[label=\arabic{section}.\arabic{subsection}.\arabic*,  align=left, labelsep=4pt, leftmargin=0pt, itemindent=57pt]

\item
Для контроля отключения всех фаз выключателя и ликвидации неполнофазного режима предусмотрена
функция защиты от неполнофазного режима (ЗНР) 

\begin{figure}[H]
\centering
\includegraphics[width=1\textwidth,height=1\textheight,keepaspectratio]{\fbpath/051.0221 ЗНР/img/ZNR.pdf}
\captionof{figure}{Функциональная схема ЗНР}\label{pict_znr:ZNR}
\end{figure}

\item
Пуск ЗНР происходит при выполнении следующих условий:

\begin{enumerate}[label=\arabic*)]
\item наличие сигнала срабатывания \textit{«ЗНФ В1 в ЗНР»} или \textit{«ЗНФ В2 в ЗНР»}; 
\item срабатывание измерительного органа, действующего по расчетному току нулевой последовательности или по отношению токов обратной и прямой последовательностей -- выбор осуществляется программным переключателем \textbf{«ИО ЗНР»};
\item отключенное положение смежного выключателя (для схем с двумя выключателями).
\end{enumerate}

\item
Наличие второго выключателя определяется программным переключателем \textbf{«Кол-во В ЗНР»}. Для
полуторной схемы также реализован контроль отключенного положения выключателя В3 конфигурируемый
вход \textit{«В3 отключен в ЗНР»}.

\item
Порог срабатывания измерительного органа по расчетному току нулевой последовательности задается
уставкой \textbf{«3I0ср ЗНР»}, по отношению токов обратной и прямой последовательностей -- \textbf{«I2/I1 ЗНР»}. Ввод в
работу защиты от неполнофазного режима осуществляется с помощью программного переключателя \textbf{«ЗНР»},
а выдержка времени срабатывания задается уставкой \textbf{«Тср ЗНР»}. 

\item
Предусмотрена возможность блокировки ЗНФ В1(2) от внешнего сигнала \textit{«Блок. ЗНФ В1(2)»} на время,
которое задается уставкой \textbf{«Тдеблок ЗНФ В1(2)»}.

\item
Оперативный ввод/вывод ЗНР производится с помощью БУ (блоков управления функциями) \textit{«Вывод ЗНР»}. 

\small
\begin{longtable}{|>{\centering\arraybackslash}m{5.3cm}|>{\centering\arraybackslash}m{3.3cm}|>{\centering\arraybackslash}m{4.2cm}|>{\centering\arraybackslash}m{1.8cm}|>{\centering\arraybackslash}m{1cm}|}
\caption{Перечень уставочных параметров ЗНР\hfill\vspace{-0.5\baselineskip}}\label{znr:tbl1}\\ 
\hline
\rowcolor{gray!30}
Параметр & Условное обозначение на схеме & Значение/ Диапазон & Единица измерения & Шаг \\ 
\hline
\endfirsthead
\caption*{\hspace{3pt}\emph{Продолжение таблицы \ref{znr:tbl1}\hfill\vspace{-0.5\baselineskip}}} \\ % сделано по ГОСТ 2.105 п.6.8.7
\hline
\rowcolor{gray!30}
Параметр & Условное обозначение на схеме & Значение/ Диапазон & Единица измерения & Шаг \\ 
%\hline
\endhead

%===m>ЗНР|-
%===m
\end{longtable}
\normalsize

\end{enumerate}